\scanpage{411}

Thus in summary
\[
(49)\quad
\left\{
\begin{array}{l}
\uline{\mathrm{Centre}}(\mathcal G^*_{\rho,\sigma}(0))=T_0\times T_\rho\times\mathbb G_m\\[5pt]
\uline{\mathrm{Centre}}\left(\underset{S\mathcal G^*_{\rho,\sigma}(0)}{\underline{\tau_{\rho,\sigma}}}\right)
=\deleted{\unsure{\mathrm{Centre}}}(T_0\times T_\rho\times\mathbb G_m)\cap\tau_{\rho,\sigma}\\[7pt]
\hspace{7em}=\left\{\left(
\underset{\substack{\in\Gamma T_0\\\simeq\Gamma\mathbb G_m}}{\underline u},
\underset{\in\Gamma T_\rho}{\underline a},
\underset{\in\Gamma\mathbb G_m}{\unsure{\lambda}}
\right)\mid u=\lambda^2,\ \deleted{\unsure{\underline a}}\inserted{\det\underline a}=\varepsilon'\unsure{\lambda^2},\ \varepsilon'^2=1\right\}
\end{array}\right.
\]

The subgroup of the center which operates \unsure{trivially} on $SG$ is thus, by (47), given by the additional restriction $u=1$, it thus identifies with \deleted{$\mu_2\times\mu_2$}
\[
(50)\quad
\left\{
\begin{array}{l}
1\times\unsure{ST'_\rho}\times\mu_2\subset\tau^{\unsure{a|0}}_{\rho,\sigma}\subset T_0\times T_\rho\times\unsure{N'_\sigma}\\[5pt]
\text{where }\unsure{ST'_\rho}\subset T_\rho\\[3pt]
\hspace{3em}\Vert\ \{\underline a\mid(\det\underline a]^2=1\}
\end{array}\right.
\]

Thus
\[
(51)\quad
\deleted{\uline{\mathrm{Centr}}(\mathcal M_{\rho,\sigma})\simeq
\left(\underset{\in T_0}{1}\times\unsure{ST'_\rho}\times\underset{\in T_\sigma}{\mu_2}\right)\times\underset{\in SG}{\mu_2}}
\]
\margin{\uline{false} the center is larger,}

The group of integral points of this \deleted{center}\inserted{group}, in the modular case, is reduced:
\[
(52)\quad
\left\{
\begin{array}{l}
\mathbb Z/6\mathbb Z\times\mathbb Z/2\mathbb Z\times\mathbb Z/2\mathbb Z\\[4pt]
\simeq\mathbb Z/3\mathbb Z\times(\mathbb Z/2\mathbb Z\times\mathbb Z/2\mathbb Z\times\mathbb Z/2\mathbb Z),
\end{array}\right.
\]
\deleted{\unsure{Group of}}
\margin{\uline{N.B.} This group is in fact contained in $\mathcal M^0_{\rho,\sigma}(\mathbb Z)$.}

product of a cyclic group of order $3$ (naturally generated by $-\rho$) with a group $(\mathbb Z/2\mathbb Z)^3$, which is none other than the group $\Sigma$ already in (\uline{IX}).

More interesting would probably be the
\margin{410}

\scanpage{412}
\mspage{507}

center of $\mathcal{M}_{\rho,\sigma}^{\circ}$, which will thus be the direct product of the center \inserted{$\mu_2$} of $SG$ (on which the \inserted{\uline{commutative}} group $\mathcal{M}_{\rho,\sigma}^{\circ}(0)\simeq\tau_{\rho,\sigma}^{\circ}$ operates \unsure{trivially}) with the subgroup of $\tau_{\rho,\sigma}^{\circ}\subset T_0\times T_\rho\times T_\sigma$ which operates \unsure{trivially} on $SG$, namely \deleted{$\unsure{\tau_{\rho,\sigma}}$}

\margin{\uline{false}}

\[
\begin{aligned}
\tau_{\rho,\sigma}^{\circ}\cap 1\times T_\rho\times\unsure{T_\sigma}
 &=\{(1,\uline{a},\uline{b})\mid\det\uline{b}=1,\,
 \deleted{(}\det\uline{a}\deleted{)^{\unsure{2}}}=1\}\\
 &\simeq ST_\rho\times ST.
\end{aligned}
\]
\margin{(and not just $(\det a)^2=1$, since we are in $\tau_{\rho,\sigma}^{\circ}$}

Thus

\unsure{We}
\deleted{\diagram{\[\begin{array}{rcl}
(53)\quad\uline{\mathrm{Centre}}(\mathcal{M}_{\rho,\sigma}^{\circ})
 &=&(1\times ST_\rho\times ST_\sigma)\times\mu_2\quad,\\
 &&\hfill\cap\qquad\\[-2pt]
 &&\hfill SG\qquad
\end{array}\]}}

\margin{\uline{false} ! Center \unsure{is} of dimension 3}

thus in the modular case we have

\deleted{\diagram{\[\begin{array}{rclccccccc}
(\unsure{54|14})\quad\uline{\mathrm{Centre}}(\mathcal{M}_{\rho,\sigma}^{\circ})(\mathbb Z)
 &\simeq&1&\times&\mathbb Z/6\mathbb Z&\times&\mathbb Z/4\mathbb Z&\times&\mathbb Z/2\mathbb Z\\
 &&&&\Vert&&\Vert&&\Vert\\
 &&&&\Gamma(ST_\rho)&&\Gamma(ST_\sigma)&&\Gamma\mu_2\subset\Gamma(SG)
\end{array}\]}}

\deleted{which is only twice as large as \deleted{$\unsure{\mathcal{M}_{\rho,\sigma}}$} the group of integral points of the center of $\mathcal{M}_{\rho,\sigma}$. \deleted{(\unsure{although the latter}}}

\deleted{\uline{Remark.} The identity component of the center of $\mathcal{M}_{\rho,\sigma}^{\circ}$ is of course $ST_\rho\times ST_\sigma$, as was suggested by the \unsure{dévissage|discussion} of the theorem on page $501'$.}

\margin{411}

\scanpage{413}
\mspage{508}

\uline{$\overline{\mathrm{XI}}$}\quad\uline{Recapitulation on the group} $\mathcal{M}_{\rho,\sigma}$.

\margin{\uline{Preliminaries}}

\textcircled{$1^\circ$} Given, over a scheme $S$, a form $G$ of $\mathrm{Gl}(2)$, hence forms $SG$, $PG$ of $\mathrm{Sl}(2)$, $\unsure{\mathrm{GP}(1)|\mathrm{PGl}(1)}$, and two sections $\widetilde\rho,\widetilde\sigma$ of $PG$, and a \unsure{1-parameter} \deleted{\unsure{unipotent|multiplicative}} additive \unsure{bordered|based} subgroup
\diagram{\begin{tikzcd}[ampersand replacement=\&,column sep=large]
\widetilde{\uline{\varepsilon}}_0:\ G_a \arrow[r,hook] \& PG
\end{tikzcd}}
\inserted{$\unsure{\varepsilon_0|G_a}$}
(which amounts to the same thing, a \deleted{\unsure{of}} regular unipotent section \unsure{of} $\varepsilon_0$ of $G$, or a regular nilpotent section $N_0$ in the form $\mathcal M$ of $\mathrm{M}_2(\mathcal O_S)$ associated to $G$) --- satisfying the conditions
\[
\begin{aligned}
(1)\quad &\widetilde\sigma^{-1}\widetilde\rho
  =\widetilde{\uline{\varepsilon}}_0(1),\qquad
  \widetilde{\uline{\varepsilon}}_0\text{ and }
  \widetilde{\uline{\varepsilon}}_1
  =\rho\bigl(\widetilde{\uline{\varepsilon}}_0\bigr)
  \bigl(=\sigma\bigl(\widetilde{\uline{\varepsilon}}_0\bigr)\bigr)
  \text{ are}\\
 &\hspace{10em}\text{independent on each fiber}\\[4pt]
(2)\quad &\mathrm{Tr}\,\sigma=0\qquad
 \text{if }\sigma\text{ is a local lifting of }\widetilde\sigma,
\end{aligned}
\]

We define \uline{canonical} liftings $\rho$, $\sigma$ of $\widetilde\rho,\widetilde\sigma$ to $G$, \deleted{\unsure{satisfy}} by the conditions
\[
(3)\quad\deleted{\unsure{\sigma\rho}}
\left\{\begin{array}{ll}
 \varepsilon_0=-\sigma^{-1}\rho
 &\text{i.e. }\varepsilon_1=-\rho\sigma^{-1}\\
 \mathrm{Tr}(\rho+\sigma)=1
 &\text{i.e. }\mathrm{Tr}\,\rho=1,
\end{array}\right.
\]
and the constructions which follow make sense.

We construct a \unsure{pinning}
\diagram{\begin{tikzcd}[ampersand replacement=\&,column sep=large]
\deleted{(\unsure{4})}\ (5)\quad\deleted{\unsure{\tau_{\rho,\sigma}}=}\ \mathrm{Gl}(2)
 \arrow[r,"\sim"] \& G
\end{tikzcd}}
by the condition that $\varepsilon_1$ take the form
\[
(6)\qquad\varepsilon_1
 =\begin{pmatrix}1&0\\1&1\end{pmatrix}
\]
whence
\[
(7)\quad\left\{\begin{array}{l}
\displaystyle\rho=\begin{pmatrix}0&1\\\unsure{v|\upsilon}&1\end{pmatrix},\qquad
 \sigma=\begin{pmatrix}0&-1\\-\unsure{v|\upsilon}&0\end{pmatrix}
 \qquad(\unsure{v|\upsilon}\in\Gamma G_m\\[6pt]
\det\rho=\det\sigma=-\unsure{v|\upsilon}
\end{array}\right.
\]
\[
(8)\qquad\varepsilon_0
 =\begin{pmatrix}1&1/\unsure{v|\upsilon}\\0&1\end{pmatrix},
\]

\margin{412}

\scanpage{414}

we consider the subgroups
\[
(9)\quad\left\{\begin{array}{ll}
 L_0=\{\uline{\varepsilon}_0(\lambda)\}
   =\left\{\begin{pmatrix}1&\lambda\\0&1\end{pmatrix}\right\}
 & (\lambda\in\Gamma G_a)\\[7pt]
 L_1=\sigma(L_0)=\rho(L_0)
   =\{\uline{\varepsilon}_1(\lambda)\}
   =\left\{\begin{pmatrix}1&0\\\lambda&1\end{pmatrix}\right\}
 & \backslash\;\unsure{\text{id}}\\[7pt]
 B'_0=\left\{\begin{pmatrix}\mu&\lambda\\0&1\end{pmatrix}
       \mid\mu\in\Gamma G_m,\,\lambda\in\Gamma G_a\right\} &\\[7pt]
 T_0=\left\{\begin{pmatrix}\mu&0\\0&1\end{pmatrix}
       \mid\mu\in\Gamma G_m\right\} &
\end{array}\right.
\]
whence
\[
(10)\qquad B'_0=T_0.L_0\qquad\text{($\frac12$-direct product)}
\]
\margin{\unsure{---}}

and we have isomorphisms
\diagram{\[
(11)\quad\left\{\begin{gathered}
\begin{tikzcd}[ampersand replacement=\&,column sep=large,row sep=large]
\uline{\varepsilon}_0:\ G_a \arrow[r,"\sim"] \& L_0
 \&\lambda\arrow[r] \&\varepsilon_0(\lambda)
      =\begin{pmatrix}1&\lambda/\unsure{v|\upsilon}\\0&1\end{pmatrix}\\
\uline{\varepsilon}_1:\ G_a \arrow[r,"\sim"] \& L_1
 \&\lambda\arrow[r] \&\varepsilon_1(\lambda)
      =\begin{pmatrix}1&0\\\lambda&1\end{pmatrix}
\end{tikzcd}\\[5pt]
\text{thus }\uline{\varepsilon}_1(\lambda)
  =\rho\unsure{(\,|\{\,}\varepsilon_0(\lambda))
  =\sigma(\varepsilon_0(\lambda)).
\end{gathered}\right.
\]}

We \deleted{set} consider
\[
(12)\quad\deleted{\unsure{\mathcal{G}_{\rho,\sigma}}}
\left\{\begin{aligned}
 T_\rho&=\text{scheme of units of }\mathcal A_\rho
  =\uline{\mathcal O}_S+\uline{\mathcal O}_S\rho\\
 &\hspace{17em}\simeq\uline{\mathcal O}_S[X]/(X^2-X-\unsure{v|\upsilon})\\
 &\deleted{\text{\unsure{of}}\ \uline{\mathcal O}_S}\\
 T_\sigma&=\text{scheme of units of }\mathcal A_\sigma
  =\uline{\mathcal O}_S+\uline{\mathcal O}_S\sigma\\
 &\hspace{17em}\simeq\uline{\mathcal O}_S[Y]/(Y^2-\unsure{v|\upsilon})
\end{aligned}\right.
\]
\diagram{\[\begin{array}{rcl}
(13)\quad &N'_\sigma&=\{g\in\Gamma G\mid g(\sigma)=\varepsilon\sigma,\ \varepsilon^2=1\}\\
 &\cap&\\[-2pt]
 &G&
\end{array}\]}
whence
\diagram{\begin{tikzcd}[ampersand replacement=\&,column sep=large]
(14)\quad 1\arrow[r] \&T_\sigma\arrow[r] \&N'_\sigma\arrow[r,"\varepsilon_\sigma"] \&\mu_2\arrow[r] \&1.
\end{tikzcd}}

We set
\[
(15)\quad\left\{\begin{aligned}
 \mathcal{G}_{\rho,\sigma}&=B'_0\times T_\rho\times T_\sigma\\[4pt]
 \mathcal{G}_{\rho,\sigma}^{*}&=B'_0\times T_\rho\times N'_\sigma.
\end{aligned}\right.
\]

\margin{413}

\scanpage{415}
\mspage{509}

\margin{$\begin{gathered}
 \mathcal M_{\rho,\sigma}[0],\ \mathcal M_{\rho,\sigma}(0)\,,\\
 \unsure{\tau_{\rho,\sigma}|\mathcal G_{\rho,\sigma}},\ \tau_{\rho,\sigma}^{\circ}\,,\\
 \mathcal M_{\rho,\sigma},\ \widetilde{\mathcal G}_{\rho,\sigma}^{+}\quad
   \deleted{\unsure{\mathcal M_{\rho,\sigma}^{\circ}}}\\
 \tau_{\rho,\sigma}^{\circ+}
\end{gathered}$}

\textcircled{$2^\circ$} We introduce
\[
\begin{aligned}
(16)\quad \mathcal M_{\rho,\sigma}[0]
 &\subset \unsure{\mathcal B_{\rho,\sigma}^{*}|\mathcal G_{\rho,\sigma}^{*}}
   =B'_0\times T_\rho\times N'_\sigma\\
 &\mathrel{\mathop{=}\limits^{\mathrm{dfn}}}
 \{(\uline u,\uline a,\uline b)\mid
   \det\uline u=\det\uline b=\varepsilon'\det\uline a,\
   \varepsilon'{}^2=1\}
\end{aligned}
\]
\[
(17)\qquad \mathcal M_{\rho,\sigma}(0)\subset\mathcal M_{\rho,\sigma}[0]
 \qquad\text{defined by }\uline u=1\quad\text{i.e.}
\]
\[
 \mathcal M_{\rho,\sigma}(0)=\mathcal M_{\rho,\sigma}[0]
   \cap(T_0\times T_\rho\times N'_\sigma)
 \qquad\text{in }\mathcal G_{\rho,\sigma}^{*}.
\]
\diagram{\begin{tikzcd}[ampersand replacement=\&,column sep=large,row sep=small]
(18)\quad L_0\arrow[r,hook] \&\mathcal M_{\rho,\sigma}[0]\\
\lambda\arrow[r,mapsto] \&(\lambda,1,1)
\end{tikzcd}}
thus we have
\[
(19)\qquad \mathcal M_{\rho,\sigma}[0]\simeq
  \mathcal M_{\rho,\sigma}\unsure{[0)|(0)}\, .\, L_0
 \qquad\text{product }\tfrac12\text{-direct}
\]

We have a homomorphism
\diagram{\begin{tikzcd}[ampersand replacement=\&,column sep=huge]
(20)\quad\unsure{\dot{\mathcal M}|\mathcal M}_{\rho,\sigma}[0]\arrow[r]
 \&G_m\times\mu_2\times\mu_2
\end{tikzcd}}
\[
 \text{denoted }\{g\mapsto\chi(g),\varepsilon_2(g),\varepsilon_3(g)
\]
for $g=(\uline u,\uline a,\uline b)$ we have
\[
\begin{aligned}
 \chi(g)&=\det\uline u=\det\uline b\\
 \varepsilon_2(g)&=\varepsilon_\sigma(b)\qquad\text{(cf. (13)(14)}\\
 \varepsilon_3(g)&=\det(\uline a)/\det\uline b\\
 &\text{i.e. }\left\{\begin{array}{l}
    \det a=\varepsilon_3\mu\\
    \text{where }\mu=\chi(g)
  \end{array}\right.
\end{aligned}
\]
trivial on $L_0$ (thus factoring through the quotient). We also set
\[
(21)\qquad\tau_{\rho,\sigma}\simeq\mathcal M_{\rho,\sigma}[0]/L_0
 \qquad\text{\unsure{where}}
\]
\margin{This notation is not at all adequate (\unsure{wrongly inspired} by $\tau_{g,\unsure{\nu|r}}$) \unsure{the sequel shows} (and what precedes already showed) that the notation $\Gamma_{\rho,\sigma}$ would be suitable, inspired by $\Gamma_{g,\nu}$ and $\Gamma_{0,3}$.}

\unsure{and} the composite
\diagram{\begin{tikzcd}[ampersand replacement=\&,column sep=large]
(22)\quad\mathcal M_{\rho,\sigma}(0)\arrow[r,hook]
 \&\mathcal M_{\rho,\sigma}[0]\arrow[r] \&\tau_{\rho,\sigma}
\end{tikzcd}}
is an \uline{isomorphism}, whence we have
\diagram{\begin{tikzcd}[ampersand replacement=\&,column sep=large]
(23)\quad\unsure{\mathcal M|\mathcal G}_{\rho,\sigma}(0)
 \simeq\unsure{\tau_{\rho,\sigma}|\tau_{\rho,\sigma}'}\arrow[r,hook]
 \&T_0\times T_\rho\times N'_\sigma
\end{tikzcd}}

\margin{414}